% Generated draft: requires mathematical and source review.
\documentclass[11pt]{article}
\usepackage[margin=25mm]{geometry}
\usepackage{fontspec}
\setmainfont{Latin Modern Roman}
\usepackage{amsmath,amssymb,mathtools,unicode-math}
\setmathfont{Latin Modern Math}
\usepackage{enumitem,needspace,xurl,hyperref}
\setlength{\parindent}{0pt}
\setlength{\parskip}{5pt}
\setlength{\emergencystretch}{3em}
\newcommand{\N}{\mathbb{N}}
\newcommand{\Z}{\mathbb{Z}}
\newcommand{\Q}{\mathbb{Q}}
\newcommand{\R}{\mathbb{R}}
\newcommand{\C}{\mathbb{C}}
\newcommand{\F}{\mathbb{F}}
\newcommand{\A}{\mathbb{A}}
\newcommand{\PP}{\mathbb{P}}
\newcommand{\E}{\mathbb{E}}
\newcommand{\eps}{\varepsilon}
\newcommand{\dd}{\,\mathrm{d}}
\begin{document}
\section*{UnsolvedMath ID: 20002328}
\textbf{Source problem code:} catalogue entry
\subsection{Definitions and mathematical statement}
Problem 12 asks for stable metric theories whose theory of beautiful pairs is not essentially separably categorical. Here separable means having a countable dense subset. A pair expands an elementary inclusion \(P\preccurlyeq M\) by predicates recording the distinguished submodel, usually distance to its sorts. Beautiful pairs are the stable-theoretic saturated and relatively saturated pairs used in the cited framework. The source's word ``essentially'' allows perturbations, and its background distinguishes exact isomorphism from approximate isomorphism. The precise perturbation system and beautiful-pair saturation convention are not included in the excerpt and must be read with the reported reference.

The supplied solution report is the following. Let \(\F_q\) be a finite field and \(V\) a vector space over \(\F_q\) of countably infinite dimension, in its vector-space language. Put \(T=\operatorname{Th}(V)\) and let \(Q=T^R\) be its randomization. For a standard atomless probability space \((\Omega,\mu)\), the random-vector sort of \(V^\Omega\) consists of measurable maps to the countable discrete \(V\), modulo almost-everywhere equality, with
\[
 d(f,g)=\mu\{\omega:f(\omega)\ne g(\omega)\}.
\]
The full randomization also includes the probability or event structure needed to interpret probabilities of formulas; the displayed random-vector sort alone is not the complete language.

The aid attributes to Hanson and Ibarlucia (2022) the published assertion that \(Q\) is \(\aleph_0\)-categorical and \(\aleph_0\)-stable in continuous logic, whereas its beautiful-pair theory \(Q_P\) is not approximately \(\aleph_0\)-categorical. In the witness described there, choose \(W\le V\) with both \(\dim W=\aleph_0\) and \(\dim(V/W)=\aleph_0\). On the product probability space, embed the one-coordinate randomizations by
\[
 f\in W^\Omega\text{ or }V^\Omega
 \quad\longmapsto\quad((\omega_1,\omega_2)\mapsto f(\omega_1))\in V^{\Omega^2}.
\]
The aid calls
\[
 (V^{\Omega^2},W^\Omega)\quad\text{a separable beautiful pair},\qquad
 (V^{\Omega^2},V^\Omega)\quad\text{a separable prime pair}.
\]
It reports that these are models of the same pair theory but admit no \(\varepsilon\)-isomorphism for \(\varepsilon<1\), in the reference's perturbation convention. The prime pair is not silently asserted here to satisfy the same saturation properties as the beautiful pair. This is a reported solution to Problem 12, with the precise pair language and meaning of \(\varepsilon\)-isomorphism left for source review.

The aid additionally states the following fully specified random-automorphism gap lemma. For every injective linear map \(h:V\to V\) with \(\dim(V/h(V))=\aleph_0\), every probability space \((X,\nu)\), and every measurable map \(\widehat g:X\to\operatorname{Aut}_{\F_q}(V)\), where the automorphism group has its pointwise-convergence Borel structure,
\[
 \sup_{v\in V}\nu\{x\in X:\widehat g(x)(v)\ne h(v)\}=1.
\]
Equivalently, for every \(\eta>0\) some \(v\) has disagreement probability greater than \(1-\eta\). This is a supremum assertion and does not require a single \(v\) attaining probability one. No independence assumption on the random automorphisms is imposed.

The original background reports that probability algebras and Hilbert spaces are \(\omega\)-stable with separably categorical beautiful pairs, and that \(L^p\) spaces have two nonisomorphic separable beautiful pairs which become isomorphic up to perturbations of the norm. The range of \(p\) and perturbation conventions are not supplied. The trailing request to build metric analogues of Hrushovski constructions is Problem 13, a separate problem, not another required conclusion of Problem 12.
\subsection{Short English statement}
The supplied literature report gives randomized finite-field vector spaces as stable, separably categorical metric structures whose beautiful-pair theory fails approximate separable categoricity. It describes two separated pair models and an additional maximal random-automorphism disagreement lemma.
\subsection{Sources}
\begin{itemize}
\item Local registry, permanent UnsolvedMath ID \(20002328\), exact public code \texttt{catalogue entry}; canonical location \texttt{aim-logic-notes.json}, \texttt{notes[103]}.
\item AIM, \emph{Questions on Model theory for metric structures}, Problem 12, page 4; supplied reference \url{https://aimath.org/WWN/continuouslogic/continuouslogic.pdf}. The source excerpt includes background on probability algebras, Hilbert spaces, and \(L^p\) spaces, followed by the distinct Problem 13. No proposer or year, original remarks, or literature field is supplied.
\item Hanson and Ibarlucia (2022), attributed by the supplied aid as a full published solution. No title, journal, theorem number, or further bibliographic data is supplied; the exact beautiful-pair and perturbation conventions require verification in that reference.
\item Supplied research attempt 1: \texttt{full\_solution}, run and registry status \texttt{solved}, rigor label \texttt{proved}, low novelty confidence; classification \texttt{SOLVED-IN-LITERATURE}, created and updated 14 August 2026. The random-automorphism lemma is an additional aid claim. The aid is machine generated, not peer reviewed, and requires independent verification before citation.
\end{itemize}
\end{document}
